\scanpage{43}
c) \inserted{\uline{Prouvé}} Modulo \uline{conjecture} d'\unsure{invariance} \unsure{birationnelle|birat. abs|birat. rel} de la ramification modérée \inserted{[si $X$ non \deleted{\unsure{pro}} projective]} [qui donne un sens au pb.] et procédant comme dans b) [\deleted{Lefschetz, résolution}\inserted{\unsure{normalisation}, Bertini...}] on est ramené à une \deleted{surface \unsure{projective}}\inserted{\uline{courbe}} lisse, -- OK.

\deleted{\unsure{et} diviseur à croisements normaux, OK.}

2) \uline{Groupe fondamental local}.

$X$ \deleted{\unsure{spectre}} schéma noeth. $\left\{\begin{array}{l}\text{\unsure{complet|complète}}\\\text{\uline{str.} \uline{local} et \uline{bon}.}\end{array}\right.$

$U$ un ouvert. A prouver

a) $\pi_1^t(U)$ de type fini

b) $\pi'_1(U)$ de présentation finie

Via résolution, et conjecture de \unsure{tameness|fameuse} (\unsure{invariance} \unsure{birationnelle|birationelle}) dans le cas a), on doit s'en tirer, en se réduisant aux énoncés analogues \deleted{\unsure{pour}} de 1)......

a) Par descente, on est ramené au cas $X$ normal, et par Lefschetz [du moins \uline{si $U=X-\text{origine}$}] au cas où $\dim X=2$. On dispose alors de Abhyankar, et cela doit marcher.....

% slanted
\margin{vérifier si c'est \uline{nécessaire}}
% archive sheet number: 42
